\section{Introduction} \label{sec:introduction}

Rapidly advancing capabilities of large language models (LLMs) raise an exciting possibility: might LLMs serve as proxies for human bidders in the development of economic mechanisms, thereby making experiments orders of magnitude cheaper to run?\footnote{An early version of this paper, not cited to protect double blind reviewing, was presented at a workshop at NeurIPS 2023.} The scientific question we address here is whether, and when, LLMs provide good human proxies in auctions. If they do, this would be practically useful: LLM simulations cost orders of magnitude less than experiments with human subjects, scale up easily, and are much faster to run. In this paper, we test whether LLMs can replicate key empirical regularities of the auctions literature. We map LLM behavior in standard auction formats against established human benchmarks from economics (and also compare with theoretical equilibria), providing a proof-of-concept for low-cost, large-scale auction experimentation using LLM bidders.

\khz{A central challenge in using LLMs for mechanism design is deciding what they should be benchmarked against. In auctions, comparing bids only to theoretical equilibrium is informative but incomplete: mechanism performance in practice depends on systematic and well-documented deviations in human behavior. Auctions therefore provide a particularly useful testbed for evaluating LLMs as human proxies. They combine sharp theoretical predictions (e.g., equilibrium bidding and dominant-strategy truth-telling) with a deep experimental and empirical literature documenting robust behavioral regularities and mistakes, including overbidding, winner’s-curse losses, and last-minute bidding under hard-close rules. While a growing literature benchmarks LLMs in strategic environments against theoretical optimality, our focus is different: we benchmark LLM behavior against human auction data and human regularities, using theory as a reference point rather than the sole criterion.}

\khz{For the core demonstration, we anchor the analysis on a single reference model, GPT-4o, since it remains a widely used, well-studied general-purpose model for strategic and economic reasoning tasks in the recent literature \citep{horton2023large, manning2024automated}.. This is a design choice for interpretability rather than a claim about model superiority. Using one canonical model helps isolate variation induced by auction format, information environment, and mechanism presentation, without conflating these effects with cross-model heterogeneity. Our objective is methodological, to test whether an out-of-the-box LLM can reproduce the comparative statics and distributional patterns observed in human auction data, and we report additional models as robustness checks.}

% As a baseline, we adopt GPT-4o as our canonical LLM. We choose a single reference model to make comparisons across mechanisms and environments interpretable and stable, and because GPT-4o remains a widely used, well-studied general-purpose model for strategic and economic reasoning tasks in the recent literature \citep{horton2023large, manning2024automated}. Our goal is not to claim that GPT-4o is the best available model at the time of writing, but rather to use a strong, widely deployed baseline to probe a methodological question: when can LLM-based simulations reproduce the \emph{patterns} found in human auction data? We additionally report results from other models to assess robustness.


In comparing with the behavior of human participants, we adopt method of moments to estimate human bidder strategies in each of different studies from experimental economics. With these estimated strategies, we populate a data-generating process that enables direct comparison of LLM bids to those inferred from decades of human laboratory experiments. Using a unified metric---the {\em scaled mean absolute deviation (SMAD)}---we study LLM behavior across first-price and second-price, sealed-bid and dynamic auction formats, and for different value environments, comparing  against experimental or empirical data from the same settings. This allows us to identify regimes in which LLM behavior resembles human behavior and where it does not, offering a methodological template for future work using LLMs as proxies for the study of human decision-making in auctions.

Our results are organized to mirror the benchmark structure in the auctions literature, moving from canonical laboratory settings to richer empirical marketplace designs.

We begin with the classic sealed-bid IPV environment and compare LLM bidding in first-price and second-price auctions to both equilibrium predictions and human laboratory benchmarks \citep{kagel1993independent, kagel2020handbook}. We find that GPT-4o captures the qualitative ranking of strategic difficulty: first-price sealed-bid auctions generate larger deviations from equilibrium than second-price auctions. At the same time, the \emph{direction} of mistakes can differ: in the second-price auction, GPT-4o exhibits systematic underbidding to a greater extent than is typically reported in the experimental literature, revealing an important divergence that would be missed by focusing only on average deviation levels.

A central theme in recent market design is that good incentives are not sufficient if agents fail to \emph{recognize} what they should do. Motivated by this perspective, we study how the same strategy-proof allocation and payment rule can induce sharply different behavior when presented through different interaction forms. We build on the notion of \emph{obvious strategy-proofness} (OSP) \citep{li2017obviously}, which formalizes when truthful play is not only dominant, but also \emph{transparent} at each decision point. In affiliated private values (APV) settings, we compare a static second-price sealed-bid auction to strategically equivalent ascending-clock implementations that instantiate the same second-price logic through sequential \texttt{stay}/\texttt{exit} decisions. We find that clock interfaces substantially reduce LLM deviations relative to sealed-bid descriptions, consistent with the idea that step-by-step threshold decisions make the dominant strategy easier to execute\citep{breitmoser2022obviousness, li2017obviously}.

We then turn to common-value auctions, where bidders observe noisy private signals about a common underlying value and face the classic winner’s-curse adverse-selection problem. Human bidders suffer systematic winner’s-curse losses that worsen with competition \citep{kagel1986winner}. We find that LLM bidders are qualitatively susceptible to the same phenomenon in first-price common-value auctions: winner profits decline as the number of bidders increases, and the winner’s curse is attenuated in second-price relative to first-price. This alignment suggests that LLM agents can reproduce nontrivial comparative statics from classic experimental designs, while still allowing us to quantify where magnitudes diverge.

Building on recent work studying how descriptions shape behavior even when incentives are unchanged, we test two interventions for second-price sealed-bid auctions: an extensive-form menu description inspired by \citet{gonczarowski2023strategyproofness}, and a clock-style description inspired by \citet{breitmoser2022obviousness}. These interventions probe whether reframing the same strategy-proof mechanism improves LLM behavior. We find a sharp asymmetry: menu-style wording yields no detectable distributional shift, while clock-framing substantially increases truth-telling and reduces dispersion, indicating that procedural reframing can be far more effective than verbal emphasis alone.

Finally, to evaluate LLMs in a richer market microstructure beyond canonical lab formats, we simulate an auction environment inspired by eBay’s proxy-bidding platform. We vary two empirically important design parameters: a hidden reserve price and a modified closing rule (soft-close, which extends the auction when bidding occurs near the end). We find that under a hard-close rule, LLM bidders concentrate decisive bid updates at the end of the auction---bid sniping---mirroring patterns documented in field data from eBay. Introducing a soft-close extension sharply reduces last-moment bidding and shifts decisive actions earlier, echoing empirical comparisons between hard-close (eBay) and soft-close (early Amazon) online auctions \citep{roth2002last}. In contrast, in our calibration low hidden reserves have limited effect on outcomes, consistent with the idea that reserves matter primarily when they bind.



We have developed a code repository to systematically run auction experiments with any number of bidders and arbitrary prompting, flexible enough to generate synthetic data for many different auction formats. For the experiments described in this paper, we ran more than 1,000 auctions with more than 5,000 GPT-4o agent participants for costs totaling less than \$400. We will release our experimental framework as an open-source benchmark for evaluating economic reasoning in  LLMs.\footnote{Repository not shared here to protect double blind review.}

The rest of the paper is arranged as follows:
Section~\ref{sec:background} reviews relevant auction benchmarks and empirical regularities. Section~\ref{sec:general_method} describes our unified metrics and simulation protocol. Section~\ref{sec:benchmark_experimental} reports results in canonical laboratory settings (sealed-bid IPV, APV/clock implementations, and framing interventions). Section~\ref{sec:winners_curse} studies common-value auctions and the winner’s curse. Section~\ref{session:ebay} presents the eBay-style marketplace benchmark. Section~\ref{sec:discussion} and  Section~\ref{sec:conclusion} conclude with implications for using LLMs as behavioral proxies and for mechanism design.